\section{Vacuum constraint correction outside a protected ball}
\label{pk:constraints}

Localized deformation of the vacuum constraints was developed by
Corvino, Corvino--Schoen, and Chru\'sciel--Delay
\cite{Corvino2000,CorvinoSchoen2006,ChruscielDelay2003}; conical localization
was established by Carlotto--Schoen \cite{CarlottoSchoen2016}. We require a
correction that fixes an inner ball and allows a decaying tail. The
explicit inverse below provides this support property with full
derivative gain. A nonlinear solve at one small base norm then gives
smooth exact constraints, even when higher data norms are large.
Section~\ref{pk2:constraints} applies the result to the rescaled beam
seed.

We use the vacuum constraint map
\[
 \Phi(g,k)=
 \bigl(R_g+(\tr_g k)^2-|k|_g^2,\,
       \nabla^j k_{ij}-\nabla_i\tr_g k\bigr).
\]
\paragraph{The exterior correction problem.}
For a nonnegative integer \(m\) and a decay exponent \(j\), put
\begin{equation}
 \|v\|_{H^m_j}
  =\|v\|_{H^m(B_2)}
   +\sup_{L\geq1} L^j
      \|v(L\,\cdot)\|_{H^m(\{1<|y|<2\})}.
 \label{pk:constraint-spaces}
\end{equation}
Componentwise versions of these norms are used for tensors.  For an
integer \(s\geq3\), let
\[
 Y^s=H^s_4(\R^3;\R\oplus T^*\R^3),\qquad
 X^s=H^{s+2}_1(\R^3;\operatorname{Sym}^2)
       \oplus H^{s+1}_2(\R^3;\operatorname{Sym}^2).
\]
The Euclidean linearization of \(\Phi\), denoted only in this subsection
by \(P\), is
\begin{equation}
 P(q,p)=
 \bigl(\partial_i\partial_jq_{ij}-\Delta q_{ii},\,
       \partial_jp_{ij}-\partial_i p_{jj}\bigr).
 \label{pk:constraint-P}
\end{equation}

\begin{proposition}[Exact correction outside a protected ball]
\label{pk:constraint-nonlinear}
Fix an integer $s\ge3$. There are constants $\epsilon_s,C_s>0$
with the following property. Let $d_B=(\delta+a_B,b_B)$ be smooth
metric and second fundamental form data on $\R^3$, with
$\delta+a_B$ positive definite and, for every multi-index $\alpha$,
\[
 |\partial^\alpha a_B(y)|\le C_\alpha\la y\ra^{-1-|\alpha|},
 \qquad
 |\partial^\alpha b_B(y)|\le C_\alpha\la y\ra^{-2-|\alpha|}.
\]
Let $\phi$ be a smooth compactly supported scalar/covector pair such that
\[
 \supp\phi\subset\{|y|\ge1\},\qquad
 \Phi(d_B)=\phi\quad\hbox{on }|y|\ge1.
\]
Define the two base sizes
\[
 \beta_s=\|a_B\|_{H^{s+2}_1}+\|b_B\|_{H^{s+1}_2},
 \qquad \eta_s=\|\phi\|_{Y^s}.
\]
If $\beta_s+\eta_s<\epsilon_s$, there is a smooth tensor pair
$v_c=(q_c,p_c)$ such that
\[
 \supp v_c\subset\{|y|\ge1\},\qquad
 \Phi(d_B+v_c)=0\quad\hbox{on }|y|\ge1,
 \qquad \|v_c\|_{X^s}\le C_s\eta_s.
\]
The corrected metric remains positive definite. The correction is smooth
across $|y|=1$ and satisfies symbol bounds of metric weight one and
second fundamental form weight two at every order. No smallness is
required of these higher bounds. No moment condition is imposed on
$\phi$, and the correction need not be compactly supported; its tail
may change the asymptotic charges. No constraint equation is asserted
for the artificially supplied datum on $|y|<1$.
\end{proposition}

The proof first constructs an exterior right inverse with full derivative
gain. Its definition on all inputs, not only supported ones, will permit
smoothness to be proved by rotations and dilations at the same base order.

\begin{lemma}[Exterior right inverse for the flat constraints]
\label{pk:constraint-linear}
There is a linear operator \(\mathcal K\), defined on every \(Y^s\) and
bounded \(Y^s\to X^s\) for each integer \(s\geq3\), such that
\begin{equation}
 \supp z\subset\{|y|\geq1\}
 \quad\Longrightarrow\quad
 P\mathcal Kz=z,\qquad
 \supp\mathcal Kz\subset\{|y|\geq1\}.
 \label{pk:constraint-right-inverse}
\end{equation}
The same construction is used at all orders.  Its conjugates by rotations
and positive dilations of the argument are smooth in the operator norm
from \(Y^s\) to \(X^s\).
\end{lemma}

\begin{proof}
We construct the inverse in three steps. Newton potentials solve the
flat constraints and supply the permitted tail. Cutting them off on the
inner ball creates a compact error whose Euclidean Killing moments
vanish automatically. An annulus-supported primitive removes that error
without changing the tail. Only this compact error has moment
conditions; the original source is unrestricted.

\paragraph{The Newton tail and its cutoff error.}
First construct an inverse without a support restriction. For
\(z=(f,F)\), let
\[
 -2\Delta u=f,\qquad \Delta W_i=F_i
\]
be the Newton potentials, and set
\begin{equation}
 q_{0,ij}=u\delta_{ij},\qquad
 p_{0,ij}=\partial_iW_j+\partial_jW_i
                 -\tfrac12\delta_{ij}\partial_kW_k.
 \label{pk:constraint-Newton}
\end{equation}
Then the scalar component of \(P(q_0,p_0)\) is \(-2\Delta u\).
Also \(p_{0,jj}=\tfrac12\partial_jW_j\), so the gradient terms in the
momentum component cancel and that component is \(\Delta W_i\).
Hence \(P(q_0,p_0)=z\).

Scaled Sobolev embedding gives
\(|z(y)|\leq C_s\la y\ra^{-4}\|z\|_{Y^s}\).  Splitting the Newton
integral into a neighborhood of the origin, a neighborhood of \(y\),
and their complement yields
\[
 |u(y)|+|W(y)|\leq C_s\la y\ra^{-1}\|z\|_{Y^s}.
\]
For completeness, on \(|v|<|y|/2\) the Newton kernel is \(O(|y|^{-1})\)
and the source has bounded total integral; on
\(|v-y|<|y|/2\), its integral against the source is \(O(|y|^{-2})\);
the remaining region has the same or a smaller bound.  Local estimates
on a fixed ball handle bounded \(y\).
On an annulus of radius \(L\), the rescaled right-hand side of the
potential equation has \(H^s\) norm \(O(L^{-2}\|z\|_{Y^s})\).
Interior elliptic estimates on a finite cover by slightly larger
rescaled balls therefore give
\[
 \|u(L\,\cdot)\|_{H^{s+2}}
       +\|W(L\,\cdot)\|_{H^{s+2}}
       \leq C_sL^{-1}\|z\|_{Y^s}.
\]
One physical derivative of \(W\) introduces an additional factor
\(L^{-1}\).  Consequently
\begin{equation}
 \|(q_0,p_0)\|_{X^s}\leq C_s\|z\|_{Y^s}.
 \label{pk:constraint-Newton-bound}
\end{equation}

Choose a smooth radial \(\chi\) which is zero on \(|y|\leq1.2\)
and one on \(|y|\geq1.8\).  With
\[
 G=P\bigl((\chi-1)(q_0,p_0)\bigr)
\]
we have \(P(\chi(q_0,p_0))=z+G\). The pair
\((\chi-1)(q_0,p_0)\) defining \(G\) is compactly supported.
Integration by parts against affine
functions and Euclidean Killing fields gives
\begin{equation}
 \int G_{\rm sc}=0,\quad \int y_aG_{\rm sc}=0,\quad
 \int G_{{\rm vec},a}=0,\quad
 \int(y_aG_{{\rm vec},b}-y_bG_{{\rm vec},a})=0.
 \label{pk:constraint-moments}
\end{equation}
Indeed two derivatives annihilate an affine scalar; one derivative
annihilates a translation, while pairing the divergence of a symmetric
tensor with a rotation gives zero.  If \(z\) is supported in
\(\{|y|\geq1\}\), then \(G=-z=0\) in \(|y|<1\), and \(G=0\) in
\(|y|>1.8\).  It is therefore supported in
\[
 A=\{1\leq|y|\leq2\}.
\]
We next give a compactly supported primitive for every pair on \(A\)
satisfying \eqref{pk:constraint-moments}.

\smallskip
\emph{A divergence primitive on the annulus.}
We use a capwise version of Bogovskii's divergence primitive
\cite{Bogovskii1980}, whose support and regularity properties are treated
by Costabel--McIntosh \cite{CostabelMcIntosh2010}. The finite mass
transfers below adapt the construction to the annulus.
Fix a unit-mass bump \(\zeta_0\) supported strictly inside \(A\).
There is a linear operator \(\mathcal D\), bounded from \(H^m\) to
\(H^{m+1}\) on inputs localized to a fixed compact set for each integer
\(m\geq0\), such that, for every scalar \(e\) supported in \(A\),
\begin{equation}
 \partial_i(\mathcal De)_i=e-\zeta_0\int e,
 \qquad \supp\mathcal De\subset A.
 \label{pk:constraint-divergence}
\end{equation}
Here is the construction and the derivative estimate.

Cover \(A\) by sufficiently narrow angular caps and take smooth
compactly supported cutoffs \(\psi_a\) whose sum is one on \(A\).
They can be chosen so that their sum is a fixed cutoff \(\eta\)
equal to one on \(A\), with support in \(1/2<|y|<3\).
For each cap choose a unit-mass bump \(\zeta_a\) near radius \(3/2\)
in its central direction.  Make the cap and bump small enough that
\[
 b\cdot\widehat v>1,\qquad |b|<2
\]
whenever \(v\in A\cap\supp\psi_a\) and \(b\in\supp\zeta_a\).
Every point on the segment from \(v\) to \(b\) then has projection
at least one on \(\widehat v\), and lies in the radius-two ball by
convexity.  The whole segment is thus contained in \(A\).

For a smooth compactly supported scalar \(e\), define
\begin{equation}
 (\mathcal D_\zeta e)(x)
  =\int e(v)(x-v)
      \int_1^\infty \zeta\bigl(v+t(x-v)\bigr)t^2\,\dd t\,\dd v.
 \label{pk:constraint-Bogovskii}
\end{equation}
Its support is contained in the union of the segments joining
\(\supp e\) to \(\supp\zeta\): if the integrand is nonzero, then
\(x=(1-t^{-1})v+t^{-1}b\) for some \(b\in\supp\zeta\).
Testing against a gradient and changing variables
\(b=v+t(x-v)\) gives
\[
 \int \mathcal D_\zeta e\cdot\nabla\varphi
   =\int e(v)\int\zeta(b)
          \bigl(\varphi(b)-\varphi(v)\bigr)\,\dd b\,\dd v.
\]
In this calculation, the substitution \(\tau=t^{-1}\) reduces the
inner integral to the integral of the derivative of
\(\varphi(v+\tau(b-v))\) for \(0\leq\tau\leq1\).
It follows that
\(\operatorname{div}\mathcal D_\zeta e=e-\zeta\int e\).

The bumps may be chosen as translates of one small bump.  Join their
centers to that of \(\zeta_0\) along paths at radius \(3/2\), and choose
the common bump small enough that every translated support remains
inside the open annulus.  If \(\gamma_a\) is such a path from the
center of \(\zeta_a\) to the center of \(\zeta_0\), then
\[
 V_a(x)=\int_0^1\gamma_a'(\tau)
                     \zeta_{\rm base}(x-\gamma_a(\tau))\,\dd\tau
\]
is smooth, compactly supported in the open annulus, and satisfies
\(\operatorname{div}V_a=\zeta_a-\zeta_0\).  Define
\[
 \mathcal De=
   \sum_a\mathcal D_{\zeta_a}(\psi_a e)
       +\sum_a V_a\int\psi_a e.
\]
For annulus-supported \(e\) this proves
\eqref{pk:constraint-divergence}.  The same formula defines
\(\mathcal D\) on arbitrary inputs after the fixed localization by
the \(\psi_a\).  Its values are always supported in a fixed compact
set, although exactness and annulus support are only needed and
asserted on the indicated inputs.

To prove the one-derivative bound, write \(r=|x-v|\) and
\(\omega=(x-v)/r\).  The kernel in
\eqref{pk:constraint-Bogovskii} is
\[
 \omega r^{-2}\int_r^\infty\zeta(v+\tau\omega)\tau^2\,\dd\tau
   =\omega r^{-2}a(v,\omega)+E(x,v),
 \qquad
 a(v,\omega)=\int_0^\infty\zeta(v+\tau\omega)\tau^2\,\dd\tau .
\]
On fixed compact sets \(a\) is smooth, and
\[
 E(x,v)
    =-(x-v)\int_0^1\zeta\bigl(v+t(x-v)\bigr)t^2\,\dd t.
\]
This remainder is jointly smooth on fixed compact sets, with bounds
controlled by smooth seminorms of the bump. The leading kernel has
degree \(-2\) in \(x-v\).  Its distributional \(x\) derivative is a
bounded multiplication term plus a principal-value kernel of degree
\(-3\), whose angular mean is zero.  To see the cancellation, integrate
the ordinary derivative over a concentric annulus: the two boundary
fluxes agree by degree \(-2\) homogeneity, so the coefficient of the
radial logarithm is zero.  Expansion of its smooth angular dependence
in spherical harmonics writes it as a summable series of mean-zero
angular convolution kernels acting on the input multiplied by smooth
coefficient functions of \(v\).  These coefficient functions have
rapidly decreasing bounds in the harmonic degree.  The classical
\(L^2\) estimate for a truncated homogeneous mean-zero singular
integral, with norm controlled by finitely many angular seminorms,
therefore bounds the differentiated kernel on \(L^2\).
The undifferentiated weak singularity is integrable, and the portions
away from the diagonal are smooth on bounded sets.  This proves
\[
 \|\mathcal D e\|_{H^1}\leq C\|e\|_{L^2}
\]
for the fixed localization.

Higher orders follow without loss by simultaneous translations.
For the unlocalized formula,
\[
 \partial_j\mathcal D_\zeta e
    =\mathcal D_\zeta(\partial_j e)
       +\mathcal D_{\partial_j\zeta}e.
\]
The same kernel estimate applies to derivatives of the bump, without
a unit-mass requirement.  Iterating this identity and differentiating
the fixed cutoffs and finite-rank terms gives
\begin{equation}
 \|\mathcal D e\|_{H^{m+1}}\leq C_m\|e\|_{H^m}
 \quad (m\geq0)
 \label{pk:constraint-divergence-bound}
\end{equation}
for the localized operator.  Approximation extends the formula and
its support assertions from smooth inputs to supported Sobolev inputs.

\smallskip
\emph{Symmetric divergence and the scalar primitive.}
Let \(F'\) be supported in \(A\), with zero translation and rotation
moments.  Set
\[
 T_{ij}=(\mathcal DF'_i)_j.
\]
Then \(\partial_jT_{ij}=F'_i\), and integration by parts gives
\[
 \int(T_{ij}-T_{ji})
       =\int(y_iF'_j-y_jF'_i)=0.
\]
Define
\begin{equation}
 D_{ijk}=\bigl(\mathcal D(T_{ij}-T_{ji})\bigr)_k,\qquad
 H_{ijk}=\tfrac12(D_{ijk}-D_{jki}+D_{kij}),\qquad
 S_{ij}=T_{ij}-\partial_kH_{ijk}.
 \label{pk:constraint-symmetric-divergence}
\end{equation}
Since \(D_{ijk}=-D_{jik}\), direct expansion gives
\[
 H_{ijk}=-H_{ikj},\qquad H_{ijk}-H_{jik}=D_{ijk}.
\]
Thus \(S\) is symmetric, and
\[
 \partial_jS_{ij}
  =F'_i-\partial_j\partial_kH_{ijk}=F'_i.
\]
All tensors remain supported in \(A\).  The two applications of
\(\mathcal D\) and the one derivative in
\eqref{pk:constraint-symmetric-divergence} show that this symmetric
divergence primitive maps \(H^m\) to \(H^{m+1}\).
The inverse trace reversal
\[
 p_{ij}=S_{ij}-\tfrac12\delta_{ij}S_{kk}
\]
satisfies \(p_{ij}-\delta_{ij}p_{kk}=S_{ij}\), and is the desired
momentum primitive.

Now let \(f'\) be supported in \(A\), with zero constant and linear
moments.  Put \(Y=\mathcal Df'\).  Then
\[
 \partial_iY_i=f',\qquad
 \int Y_a=-\int y_af'=0.
\]
Set
\[
 M_{ab}=\int(y_aY_b-y_bY_a),\qquad
 C_{ab}=-\tfrac12M_{ab},\qquad
 \widehat Y_i=Y_i+\partial_j(C_{ij}\zeta_0).
\]
The skew-symmetry of \(C\) implies
\[
 \partial_i\widehat Y_i=f',\qquad \int\widehat Y_i=0.
\]
Moreover,
\[
 \int\bigl(y_a\partial_j(C_{bj}\zeta_0)
              -y_b\partial_j(C_{aj}\zeta_0)\bigr)=2C_{ab},
\]
so every rotation moment of \(\widehat Y\) vanishes.
Apply the symmetric divergence primitive to \(\widehat Y\), obtaining
\(S\) with \(\partial_jS_{ij}=\widehat Y_i\), and put
\(q_{ij}=S_{ij}-\tfrac12\delta_{ij}S_{kk}\).  Then
\[
 \partial_i\partial_jq_{ij}-\Delta q_{ii}
    =\partial_i\partial_jS_{ij}=f'.
\]
This gains two derivatives and preserves support in \(A\).

Together these constructions give an operator \(\mathcal B\) satisfying
\begin{equation}
 P\mathcal B(f',F')=(f',F'),\qquad
 \|\mathcal B(f',F')\|_{H^{s+2}\oplus H^{s+1}}
    \leq C_s\|(f',F')\|_{H^s\oplus H^s}
 \label{pk:constraint-compact-primitive}
\end{equation}
on the supported moment-free pairs, with output supported in \(A\).
They also define a bounded operator on arbitrary compact inputs:
use the fixed localization in every occurrence of \(\mathcal D\),
and symmetrize the final tensors.  On moment-free annulus-supported
inputs this final symmetrization changes nothing.  Off that subspace
no exactness is asserted or required.

We may now set
\begin{equation}
 \mathcal Kz=\chi(q_0,p_0)-\mathcal BG.
 \label{pk:constraint-K}
\end{equation}
The estimates already proved give its boundedness on all \(Y^s\).
For exterior-supported inputs, both the support assertion and
\(P\mathcal Kz=z\) follow from
\eqref{pk:constraint-compact-primitive} and \(P(\chi(q_0,p_0))=z+G\).

Finally, all the operators have smooth conjugates under the stated
group actions.  For the argument dilation
\((T_\lambda e)(y)=e(\lambda y)\), substitution in
\eqref{pk:constraint-Bogovskii} gives
\begin{equation}
 T_\lambda\mathcal D_\zeta T_\lambda^{-1}
        =\lambda\mathcal D_{\zeta_\lambda},\qquad
 \zeta_\lambda(y)=\lambda^3\zeta(\lambda y).
 \label{pk:constraint-conjugation}
\end{equation}
The transformed cutoffs, path fields, and moment functionals likewise
depend smoothly on \(\lambda\).  Rotations transform the bumps and
cutoffs smoothly and multiply tensor components by rotation matrices.
The kernel bounds above are controlled by finitely many smooth
seminorms on fixed compact sets, so Taylor's formula for these
parameter-dependent coefficients proves smoothness in the operator
norm with the full derivative gain.  Newton inversion conjugates by
the factor \(\lambda^2\), and its first derivative by the corresponding
factor \(\lambda\); it is equivariant under rotations.
Composition in \eqref{pk:constraint-K} proves the final assertion.
\end{proof}

\begin{proof}[Proof of \Cref{pk:constraint-nonlinear}]
The same fixed point will supply both the base estimate and the
all-order regularity.

\paragraph{A contraction at one base order.}
Use the operator $\mathcal K$ from \Cref{pk:constraint-linear}.
Define
\begin{equation}
 N(q,p)=\Phi(d_B+(q,p))-\Phi(d_B)-P(q,p).
 \label{pk:constraint-nonlinearity}
\end{equation}
For small \(\beta_s\) and on a sufficiently small \(X^s\) ball this
is a smooth map to \(Y^s\).  We record the weights and derivatives
to justify this assertion.  Metric perturbations have weight one,
their first and second derivatives have weights two and three,
and second fundamental form perturbations and their first derivatives have weights
two and three.  After subtraction of \(P\), the scalar terms are
linear combinations, with smooth bounded inverse-metric factors,
of terms with the budgets
\[
 a\,\partial^2a:\ 1+3=4,\qquad
 (\partial a)^2:\ 2+2=4,\qquad
 b^2:\ 2+2=4.
\]
The corresponding momentum budgets are
\[
 a\,\partial b:\ 1+3=4,\qquad
 (\partial a)b:\ 2+2=4.
\]
Here \(a,b\) may be background or correction factors, with the
background-only terms canceled in
\eqref{pk:constraint-nonlinearity}.  In particular, a coefficient in
front of a highest derivative differs from its Euclidean value by
a weight-one factor.  Applying the Sobolev product estimates on
the unit ball and on the rescaled annuli, using \(s\geq3\), yields
\begin{equation}
 \|N(v)-N(w)\|_{Y^s}
   \leq C_s(\beta_s+R)\|v-w\|_{X^s},
 \qquad \|v\|_{X^s},\|w\|_{X^s}\leq R.
 \label{pk:constraint-Lipschitz}
\end{equation}
The same products after differentiation with respect to \(v\) prove
smoothness.  Inverse-metric factors are smooth because the metrics
in this ball remain uniformly positive.  Notice that the assertion
concerns \(N\); the uncanceled Euclidean leading terms of \(P\) have
weight three on a general element of \(X^s\).

Let \(C_{\mathcal K,s}\) be the norm of \(\mathcal K:Y^s\to X^s\).
For sufficiently small \(\beta_s,\eta_s\), the equation
\begin{equation}
 z_c=-\phi-N(\mathcal Kz_c)
 \label{pk:constraint-fixed-point}
\end{equation}
is a contraction on \(\|z_c\|_{Y^s}\leq2\eta_s\).  Indeed take
\(R=2C_{\mathcal K,s}\eta_s\) in
\eqref{pk:constraint-Lipschitz} and require
\(C_sC_{\mathcal K,s}(\beta_s+R)<1/2\).
Starting the iteration at zero preserves support in \(\{|y|\geq1\}\):
\(\phi\) has that support,
\Cref{pk:constraint-linear} gives it for every correction, and
\(N(v)=0\) pointwise wherever \(v\) vanishes on an open set.
The limit therefore has the same support.  With \(v_c=\mathcal Kz_c\),
\begin{equation}
 \|v_c\|_{X^s}\leq2C_{\mathcal K,s}\eta_s,\qquad
 \Phi(d_B+v_c)=0\quad\hbox{on } |y|\geq1.
 \label{pk:constraint-correction-bound}
\end{equation}
The equation holds there because \(Pv_c=z_c\) and
\(\Phi(d_B)=\phi\) there.  On the inner side \(v_c=0\).

\paragraph{Higher smoothness from rotations and dilations.}
It remains to establish smoothness at every order, without imposing
smallness of all high Sobolev norms.  We do this at the single base
order \(s\).  Let \(U_g\) and \(V_g\) be the actions of a rotation and
a positive argument dilation on tensor pairs and on scalar/covector
pairs respectively.  Rotations act on components as well as arguments;
dilations here act on the argument alone.  These actions are uniformly
bounded on \(X^s,Y^s\) for \(g\) near the identity.  Set
\[
 \mathcal K_g=U_g\mathcal K V_g^{-1},\qquad
 \phi_g=V_g\phi,\qquad
 N_g(w)=V_gN(U_g^{-1}w).
\]
The operator family \(\mathcal K_g\) is norm-smooth by
\Cref{pk:constraint-linear}.  All background coefficients have the assumed smooth symbol bounds,
and the source is smooth and compactly supported.  Thus \(g\mapsto\phi_g\) is smooth in \(Y^s\).
Writing \(N_g\) using the weighted products displayed above shows
that it is jointly smooth in \(g,w\): parameter derivatives fall on
the smooth background coefficients or introduce the explicit scale
and component factors in a transformed derivative.  They do not
differentiate the unknown \(w\) beyond the derivatives already
allowed in \(X^s\).

Consider the equation on the full space \(X^s\),
\begin{equation}
 w+\mathcal K_g\phi_g+\mathcal K_gN_g(w)=0.
 \label{pk:constraint-parameter-equation}
\end{equation}
At \(g\) equal to the identity its solution is \(v_c\).
The derivative in \(w\) is
\(\Id+\mathcal K_gD N_g(w)\), whose inverse is a convergent Neumann
series on a sufficiently small common ball by
\eqref{pk:constraint-Lipschitz}.  The implicit function theorem
therefore gives a smooth solution of
\eqref{pk:constraint-parameter-equation} in the group parameters.

We specify the uniqueness argument because argument dilations move
the support boundary.  First choose a fixed small ball on which
the contraction estimates, and their nearby transformed versions,
hold.  Taking the defect sufficiently small puts \(v_c\) strictly
inside a smaller ball.  The uniform boundedness of \(U_g\) then puts
every \(U_gv_c\), for \(g\) near the identity, in the common contraction
ball.  It satisfies \eqref{pk:constraint-parameter-equation} by
conjugation.  Uniqueness in that ball identifies it with the smooth
implicit-function branch.  This uses the all-input definition of
\(\mathcal K\), so no fixed support subspace has to be invariant
under dilation.  Nor does it presuppose norm continuity of the
group orbit of the unknown solution.

Distributional differentiation of this smooth orbit shows that all
iterated rotations
\(\Omega_{ij}=y_i\partial_j-y_j\partial_i\) and dilations
\(E=y\cdot\nabla\), applied to \(v_c\), lie in \(X^s\).
The rotation of tensor components contributes only constant
zeroth-order matrices, which may be subtracted in identifying these
derivatives.
On every rescaled exterior annulus these fields span the scaled
coordinate derivatives.  More explicitly, with \(r=|y|\),
\[
 r^2\partial_i=y_iE+\sum_j y_j\Omega_{ji}.
\]
Their iterates therefore control every scaled spatial derivative.
The same identity is nondegenerate on a neighborhood of \(|y|=1\).
It follows that the correction is smooth across its support
boundary as well as on the exterior, and has symbol bounds of all
orders with metric weight one and second fundamental form weight two.
These higher constants may depend on the higher symbol bounds of
$d_B$ and on the smooth source $\phi$; their smallness is not required.
The small base norm also keeps the corrected metric positive definite.
\end{proof}

\paragraph{Application to the beam seed.}
The rescaling, artificial interior extension, and verification of the
source hypotheses are carried out in Section~\ref{pk2:constraints},
using \eqref{pk:constraint-rescaling}--\eqref{pk:constraint-background}.
That application retains the complete correction tail and proves
approximation in any prescribed finite list of the original data
seminorms. Higher symbol bounds ensure smooth membership of each
corrected datum; they need not be uniformly small or give convergence
at all orders for one fixed frequency exponent.
